\subsection{Overview}
We model route generation as a decision--execution cascade with an explicit \emph{corridor commitment}.
Let $y$ denote a graph path (node sequence) and $\tau$ denote a continuous trajectory realization (optional).
Our factorization is:
\begin{equation}
  p_\theta(y \mid c) = \sum_{z} p_\theta(z \mid c)\, p(y \mid z, c),
\end{equation}
where $z$ is a \emph{sparse waypoint commitment} on the road graph.
The key design choice is to make the decision stage \emph{short-horizon} (few steps) to avoid the compounding error of long node-level autoregressive generation.

\subsection{Decision: waypoint-level autoregressive commitment}
\paragraph{Ground-truth waypoint extraction.}
Given a ground-truth graph path $y=(v_0,\dots,v_L)$, we derive a fixed-length waypoint sequence
$z=(w_0,\dots,w_M)$ with $w_0=v_0$ and $w_M=v_L$.
We use a fixed-$K$ variant of the Ramer--Douglas--Peucker procedure in graph coordinate space: iteratively select the intermediate nodes with the largest perpendicular deviation from the current polyline until $M$ reaches the desired number of waypoints.\footnote{This is a supervision heuristic to stabilize training under limited data; learning waypoints end-to-end is left to future work.}

\paragraph{Waypoint bins as a compact output space.}
Directly predicting a waypoint as a node ID is a classification problem over hundreds of thousands of nodes.
We instead discretize the grid into $B\times B$ waypoint bins (e.g., $B=32$) and train an autoregressive classifier over $B^2$ classes.
At step $m$, the model predicts the next waypoint bin conditioned on the previous waypoint, the destination, and context features from $t_0$ and road tier.
This reduces the autoregressive horizon to $M\!\approx\!3$--$5$ steps, making exposure bias and compounding error much less severe than node-level AR over hundreds of edges.

\subsection{Execution: A* connection on the road graph}
Given a sampled waypoint sequence $\hat{z}=(\hat{w}_0,\dots,\hat{w}_M)$, we execute it into a full corridor path by graph search.
For each consecutive pair $(\hat{w}_{m},\hat{w}_{m+1})$, we run A* on $G$ with edge-length costs and concatenate the resulting segments into a complete path $\hat{y}$.
If any segment is unreachable, we mark the sample as a failure; this motivates reachability-aware decoding as a key next step.
Because the executor enforces that the final path passes through $\hat{z}$ and stays on $G$, it is \emph{commitment-preserving} by construction.

\subsection{Optional continuous executor (extension)}
The graph path $\hat{y}$ can be converted back to a continuous polyline and refined with a continuous generator (diffusion/flow) to model within-corridor variation (lane choice, smoothing, local detours).
We treat this as an optional extension: the current paper focuses on recovering corridor-level structure with explicit commitments and reliable execution on the road graph.
